% ============================================
% Johns Hopkins Letter Template
% Assistant Professor Cover Letter – Colorado College
% ============================================

\documentclass[11pt,letterpaper]{letter}

\usepackage{graphicx}
\usepackage{eso-pic}
\usepackage{geometry}
\usepackage{microtype}
\usepackage[hidelinks]{hyperref}

% ----- Page Setup -----
\geometry{left=25mm,right=25mm,top=25mm,bottom=25mm}

% ----- Logo Placement (first page only) -----
\newcommand{\PutJHULogoFG}[1]{%
  \AddToShipoutPictureFG*{%
    \AtPageUpperLeft{%
      \hspace{10mm}%
      \raisebox{-38mm}{%
        \includegraphics[width=6cm]{#1}%
      }%
    }%
  }%
}

% ----- Sender Information -----
\signature{Decory Edwards}
\address{Department of Economics \\ Johns Hopkins University \\ Baltimore, MD 21218 \\ \texttt{dedwar65@jh.edu}}

\begin{document}
\PutJHULogoFG{Images/jhulogo.png}

\begin{letter}{Search Committee \\
Department of Economics and Business \\
Colorado College \\
14 E. Cache La Poudre St. \\
Colorado Springs, CO 80903 \\ USA}

\opening{Dear Members of the Search Committee,}

I am writing to express my interest in the tenure-track Assistant Professor position in the Department of Economics and Business at Colorado College, beginning Fall 2026. I am a Ph.D. candidate in Economics at Johns Hopkins University, advised by Professors Christopher D.~Carroll, Nicholas W.~Papageorge, and Stelios Fourakis, and I expect to receive my degree in May 2026. My research fields are macroeconomics, wealth and income dynamics, and computational economics.

My research agenda focuses on understanding the determinants of wealth inequality and the mechanisms through which heterogeneity in returns to wealth shapes aggregate outcomes. In my job-market paper, I estimate the degree of heterogeneity in individual portfolio returns using micro-data from the Survey of Consumer Finances and simulate a structural model in which households face idiosyncratic return risk. I show that allowing for persistent heterogeneity in returns substantially improves the model’s ability to match the empirical distribution of wealth, offering a tractable way to connect micro-level return dispersion to macroeconomic inequality.

In a second project, I use the Health and Retirement Study to examine the relationship between interpersonal trust and portfolio performance. Building on the literature that finds a hump-shaped relationship between trust and income, I show that a similar relationship holds for individual returns to wealth measured in the HRS data, highlighting a behavioral channel through which social attitudes shape saving and investment outcomes. This work also provides new estimates of the persistent component of returns to net worth, which motivate the calibration of heterogeneity in my structural model.

Colorado College’s distinctive Block Plan, immersive pedagogy, and commitment to an inclusive and intellectually vibrant liberal arts environment make this position especially appealing. I am enthusiastic about the opportunity to teach macroeconomics within a structure that allows deep engagement with students and supports innovative formats such as field-based learning, intensive group projects, and community-engaged work.The position’s emphasis on teaching international finance, money and banking, principles of macroeconomics, and intermediate macroeconomic theory aligns directly with my background and training.

\textbf{Teaching.} My teaching experience centers on undergraduate macroeconomics and an upper-level macro policy seminar, where I have led weekly discussion sections and held regular office hours for classes ranging from 16 to 40 students. My approach is student-centered: I aim to help students “convince themselves” that they have moved from confusion to understanding by holding structured office-hour coaching that builds trust and confidence. At Colorado College, I am prepared to teach \textit{Principles of Macroeconomics}, \textit{Intermediate Macroeconomic Theory}, and \textit{Money and Banking}, as well as to supervise senior theses and advise students across block-based research and independent projects.

I believe I would be a strong fit because my computational and empirical toolkit allows me to (i) deliver rigorous, data-driven instruction across macroeconomics and quantitative methods, (ii) integrate reproducible workflows and real-world economic applications into my courses, and (iii) mentor students effectively in a liberal-arts environment that values inclusive pedagogy, student engagement, and high-impact teaching. I am excited about the opportunity to contribute to Colorado College’s teaching, research, and service missions within a community committed to equity and belonging.

I would be honored to join the Department of Economics and Business at Colorado College. Thank you for your consideration; I welcome the opportunity to discuss my fit with your department at your convenience.

\closing{Sincerely,}

\end{letter}
\end{document}
